% page 597 du fac-simile (image p-596 du scan)
% bloc 162.6 x 242.0 mm ; marge gauche 39.4 mm
\pgc{17.780mm}{0.000mm}{
\l{\cel{49}589.}
\l{}
\l{\sou{traktoro}. (\sou{.tekn.}) Vehilo automobila di qua la roti esas pro-}
\l{\cel{3}vizita ye organi por igar lu adherar a sulo mem tre ne-}
\l{\cel{3}plana (roti qui inter-relatas per bendo metala quaze-tapisa}
\l{\cel{3}qua, lor la funciono, aspektas quale raupo movanta - o roti}
\l{\cel{3}angul-pecoza) e qua posedas junto-sistemo por remorko (di}
\l{\cel{3}plugilo, di falcho-mashini, e c.) -\cel{2}DEFIRS.}
\l{}
\l{\sou{traktrico}. (\sou{geom})\cel{2}Kurvo quan trasas kordo quan onu tir-esforcas}
\l{\cel{3}- altra-vorto : di qua la tangento egalesas lineo konstanta.}
\l{\cel{3}- DEFIS.}
\l{}
\l{\sou{tramo.}\cel{2}(\sou{tekn.)}\cel{2}Longa stango de fero o, maxim-multa-kaze, de}
\l{\cel{3}stalo, en qua existas kanelo por recevar la parto salianta}
\l{\cel{3}di la felgo di veturo-roto, e qua, instalite en sulo, ne}
\l{\cel{3}salias de la surfaco di olca. (opoze a \textquotedbl{}relo\textquotedbl{}). - DEFIS.}
\l{}
\l{\sou{tramplar}. (\sou{netrans.}) Movar, agitar, sua pedi, sen avancar.}
\l{\cel{3}- DE.}
\l{}
\l{\sou{trampolino}. (\sou{tekn.)}\cel{3}Planajo pento-forma, ek planko instalita}
\l{\cel{3}tale ke olu esas elastika, de-sur qua la saltanti ri-saltas,}
\l{\cel{3}tale aquirinte springo plu granda. - DFIRS.}
\l{}
\l{\sou{tranar}. (\sou{trans.}) I. Tirar dop su (ulo, ulu) malgre la rezisto}
\l{\cel{3}de la pezo o de la froto di parto (de lu) sur sulo. - II.}
\l{\cel{3}Koaktar (ulu) pri kun-venar. - EF.}
\l{}
\l{\sou{tran}car. (\sou{patol.}) Nekonciar, dum epilepsio, dum hipnoto, e c.}
\l{\cel{3}DEFIRS.}
\l{}
\l{\sou{trancendanta}. (\sou{filoz.})\cel{2}La filozofio qua traktas exkluzive}
\l{\cel{3}la konoco aprioria, la konoco pure racionala. (En la sistemo}
\l{\cel{3}da Immanuel Kant, spaco e tempo esas koncepti trancendanta).-}
\l{\cel{3}DEFIRS.}
\l{}
\l{\sou{tranchar}. (\sou{trans.)}\cel{3}Dividar par-separe (korpo) per korpo plu}
\l{\cel{3}harda e taliita tale ke adminime un ek lua aristi esas akuta.}
\l{\cel{3}(ex.: tranchar la kapo per gilotino, per hakilo-frapo). EFIS.}
\l{}
\l{\sou{trancheo}. (\sou{tekn.}) I. Exkavuro facita en sulo ula-longe (por}
\l{\cel{3}kanalo, kanaleto, voyo, relvoyo, e c.) -II. (\sou{militarto}).}
\l{\cel{3}Fosato, shirme da qua onu povas proximigar su a fortifikajo,}
\l{\cel{3}a fortreso, +asiejata (siejata). - DEFR.}
\l{}
\l{\sou{tranquila}.\cel{2}Qua ne esas agitema, bruisema. - EFIS.}
\l{}
\l{\sou{trans}. Plu fore kam. - L.}
\l{}
\l{\sou{transaktar}. (\sou{netrans., pri}). Abutar ad interkonsento per cedi}
\l{\cel{3}reciproka pri la litijo-punto. - EFIS.}
\l{}
\l{\sou{transbordar}. (\sou{trans.)}\cel{2}Igar (vari, pasajanti) pasar de ca a ta}
\l{\cel{3}navo - o de loko ube la relvoyo esas interruptita, a la loko}
\l{\cel{3}ube lu esas itere uzebla. - FS.}
}
